% Figure: the cell as a control loop. Sensor -> Controller -> Actuator
% -> Output -> Feedback. The canonical block diagram of classical
% control engineering, redrawn with cellular components named at each
% block.
\begin{figure}[htbp]
\centering
\begin{tikzpicture}[
  block/.style={draw=engblack, thick, rectangle, minimum width=2.6cm,
    minimum height=1.1cm, align=center, font=\small, fill=white},
  sense/.style={block, fill=engblue!10},
  ctrl/.style={block, fill=engamber!15},
  act/.style={block, fill=engred!10},
  sum/.style={draw=engblack, thick, circle, minimum size=0.6cm,
    inner sep=0pt, font=\small},
  arr/.style={-{Latex[length=2.2mm]}, thick, draw=engblack},
  lbl/.style={font=\scriptsize\itshape, align=center},
]
% Top row: forward path
\node[sum] (sum) at (0,0) {$\Sigma$};
\node[ctrl] (ctrl) at (3,0) {Controller \\ \scriptsize (transcription \\ \scriptsize factor cascade)};
\node[act] (act) at (6.7,0) {Actuator \\ \scriptsize (enzyme; channel; \\ \scriptsize pump)};
\node[block, fill=englime!15] (plant) at (10.3,0) {Process \\ \scriptsize (the cell's \\ \scriptsize chemistry)};

% Output
\node[block, fill=englime!20] (out) at (10.3,-2.4) {Output \\ \scriptsize (concentration; \\ \scriptsize voltage; rate)};

% Sensor
\node[sense] (sens) at (3,-2.4) {Sensor \\ \scriptsize (receptor; \\ \scriptsize metabolic sensor)};

% Setpoint and disturbance
\node[lbl, anchor=east] (sp) at (-1.4,0) {setpoint \\ (genetic / \\ homeostatic)};
\node[lbl, anchor=south] (dist) at (10.3,1.4) {disturbance \\ (stress; toxin; \\ ligand)};

% Arrows forward
\draw[arr] (sp) -- (sum);
\draw[arr] (sum) -- (ctrl) node[midway, above, font=\scriptsize] {error};
\draw[arr] (ctrl) -- (act);
\draw[arr] (act) -- (plant);
\draw[arr] (dist) -- (plant);
\draw[arr] (plant) -- (out) node[midway, right, font=\scriptsize] {measured};

% Feedback path
\draw[arr] (out) -- (sens);
\draw[arr] (sens) -| node[near end, left, font=\scriptsize] {$-$} (sum.south);

\end{tikzpicture}
\caption[The cell as a control loop]{The cell as a control loop. A
setpoint (encoded genetically or homeostatically) is compared against
a sensor reading; the error drives a controller (a transcription
factor cascade, a second-messenger cascade, or a covalent modification)
that drives an actuator (an enzyme, an ion channel, a pump) that acts
on the process (the cell's internal chemistry). The output is measured
by a sensor whose reading closes the loop. Each block has a named
cellular implementation, and the cell's behaviour at the loop level is
governed by the same equations as the engineering analogue.}
\label{fig:vol06:ch01:control-loop}
\end{figure}
